\documentclass[a4paper,12pt]{article}
\usepackage[a4paper,margin=1in]{geometry}
\usepackage{titlesec}
\usepackage{graphicx}
\usepackage{enumitem}
\usepackage{hyperref}
\usepackage{fancyhdr}
\usepackage{amsmath}
\usepackage{xcolor}
\usepackage{microtype}

\pagestyle{fancy}
\fancyhf{}
\rhead{Pathfinder}
\lhead{User Manual}
\cfoot{\thepage}

\setlist[itemize]{noitemsep,topsep=2pt,leftmargin=1.2cm}

\title{\vspace{-2cm}\textbf{Pathfinder}\\\large User Documentation and Technical Summary}
\author{Developed in C++}
\date{}

\begin{document}
\maketitle
\vspace{-0.4cm}
\hrule
\vspace{0.4cm}

\section*{1. Overview}
\textbf{Pathfinder} is a C++ simulation of an emergency response network using a grid-based graph model of a city.  
Each grid node represents an intersection, house, hospital, or ambulance station. The system uses \textbf{Dijkstra’s shortest path algorithm} to route ambulances and find hospitals during emergencies, while dynamically responding to disasters such as earthquakes and floods that damage parts of the grid.

\section*{2. Menu Options}
\begin{itemize}
    \item \textbf{ (1)Add Ambulance:} Assign a node as an ambulance station.
    \item \textbf{ (2)Add House:} Mark a node as a house.
    \item \textbf{ (3)Add Hospital:} Assign a node as a hospital.
    \item \textbf{ (4)Trigger Emergency:} Simulate an incident and trace optimal paths.
    \item \textbf{ (5)Simulate Earthquake:} Damage nearby nodes randomly.
    \item \textbf{ (6)Simulate Rain:} Flood nodes with poor drainage.
    \item \textbf{ (7)Show Damaged Nodes:} Display damaged or flooded nodes.
    \item \textbf{ (8)Exit:} Quit the simulation.
\end{itemize}

\section*{3. Compilation and Execution}
\textbf{To Compile:} \begin{verbatim} g++ main.cpp GridGraph.cpp -o Pathfinder \end{verbatim}
\textbf{To Run:} \begin{verbatim} ./Pathfinder \end{verbatim}

\section*{4. Example Run}
\begin{verbatim}
==========================================
     EMERGENCY RESPONSE SIMULATION
==========================================
Enter grid size: 4

4x4 grid created successfully!
Default: Node 0 = Ambulance | Node 15 = Hospital

Enter your choice: 2
Enter node ID to assign as House : 5
 House assigned to node 5

Enter your choice: 4
Enter house node ID for emergency: 5

 Emergency triggered at node 5!
Finding shortest path from ambulance to emergency house...
Path: AmbulanceStation -> Intersection_4 -> House_5
Distance: 9
Finding nearest hospital from the emergency house...
Path: House_5 -> Intersection_9 -> CityHospital
Distance: 12
\end{verbatim}

\section*{5. Tips for Effective Simulation}
\begin{itemize}
    \item \textbf{Setup before running:} Ensure that at least one \textbf{ambulance} and one \textbf{hospital} node exist before triggering an emergency to guarantee valid rescue routes.
    \item \textbf{Experiment safely:} Start with small grids such as $4\times4$ or $5\times5$ to clearly observe pathfinding and disaster effects.
    \item \textbf{Interpret disasters:} Damaged or flooded nodes (type = \texttt{DISASTER}) are permanently inaccessible until a new simulation run.
    \item \textbf{Randomization note:} Each simulation run uses random edge weights and drainage values, so paths and damages differ every time.
    \item \textbf{Performance tip:} Larger grids increase route-finding complexity; use moderate sizes for smoother interaction.
    \item \textbf{Debugging help:} If routes seem incomplete, check whether key nodes (house, ambulance, hospital) were affected by disasters.
\end{itemize}


\section*{6. Technical Summary}
\begin{itemize}
    \item \textbf{Language \& Tools:} Implemented in C++ using STL containers, priority queues, and random number generation.
    \item \textbf{Core Algorithm:} Dijkstra’s algorithm for shortest path computation.
    \item \textbf{Structure Overview:}
    \begin{itemize}
        \item \texttt{NodeStructures.h} – Defines node attributes (ID, type, name, drainage).
        \item \texttt{GridGraph.h /.cpp} – Implements grid construction, pathfinding, and disasters.
        \item \texttt{main.cpp} – Provides the interactive menu-driven simulation.
    \end{itemize}
    \item \textbf{Disaster Handling:}
    \begin{itemize}
        \item Earthquake: Damages nodes within a radius of the epicenter.
        \item Rain: Floods low-lying nodes based on random drainage levels.
    \end{itemize}
\end{itemize}

\vfill
\noindent\textit{Developed for educational purposes to demonstrate graph algorithms, routing, and disaster simulation.}

\end{document}